\clearpage
\subsection{Trazabilidad entre objetivos, método y productos}
\label{sec:trazabilidad-objetivos}
Los cuatro objetivos se articulan en una secuencia: configurar la matriz, aplicarla al diseño inicial, formular una propuesta que responda al diagnóstico y validar sus efectos. Simular casos no sustituye la aplicación de los indicadores ni resuelve automáticamente el desempeño ambiental de los materiales.

\begin{table}[H]\centering
\caption{Correspondencia entre objetivos específicos y productos verificables}
{\small\singlespacing\renewcommand{\arraystretch}{1.35}\setlength{\tabcolsep}{4pt}
\begin{tabular}{>{\raggedright\arraybackslash}p{2.6cm}>{\raggedright\arraybackslash}p{5.6cm}>{\raggedright\arraybackslash}p{5.8cm}}\toprule
Objetivo & Método & Producto y criterio de cierre \\ \midrule
OE1. Configurar la matriz & Selección, adaptación y documentación de indicadores de los cinco marcos. & Capítulo 3: matriz con fuente, versión, unidad, cálculo, referencia y regla de decisión. \\
OE2. Evaluar el CB & Aplicación de cada indicador al modelo y documentación del diseño inicial. & Capítulo 4: resultados, brechas y no evaluables por indicador, sin confundir línea base y certificación. \\
OE3. Desarrollar la propuesta & Medidas pasivas y materiales seleccionados a partir de las brechas, con viabilidad documentada. & Capítulo 5: escenarios trazables y especificaciones técnicas y ambientales justificadas. \\
OE4. Validar la incidencia & Comparación energética y de confort en condiciones equivalentes; revisión de límites e incertidumbres. & Capítulo 5: diferencias de energía y de horas aceptables por zona; síntesis de alcance en conclusiones. \\
\bottomrule\end{tabular}}
\fuente{Elaboración propia a partir de los objetivos específicos de la investigación.}
\end{table}
